\documentclass[11pt]{article}
\usepackage[margin=1in]{geometry}
\usepackage[T1]{fontenc}
\usepackage{lmodern,microtype}
\usepackage[colorlinks=true,allcolors=black]{hyperref}
\pagestyle{empty}
\begin{document}
\noindent Mathias Rodr\'iguez Castro\\
Facultad de Ingenier\'ia, Universidad de la Rep\'ublica\\
Montevideo, Uruguay\\
\href{mailto:mathiasr@fing.edu.uy}{mathiasr@fing.edu.uy}

\bigskip
\noindent To the Editors of \emph{Operations Research Forum}

\bigskip
\noindent Dear Editors,

\medskip
Please consider ``Accuracy-Aware Row Scaling for Generated Mixed-Integer Programs:
A Computational and Software Study'' as an Original Research Article in
\emph{Operations Research Forum}.

Operational models generated in software require a connection between numerical
feasibility and accuracy in the application's original units. Row scaling preserves
the mathematical model in exact arithmetic but changes that connection. The paper
develops a reproducible workflow that makes row-level residual budgets explicit
inputs to factor selection, export checks and independent returned-point verification.

The first contribution is accuracy-aware model generation and verification, supported
by admissibility conditions, incompatibility checks, integer-rounding allowances
and a selector checking finite-export error reserves. The second is a computational study that separates
conditioning, subsequent equilibration, optimizer effort and complete-process cost.
The mathematical results support these implementation decisions and make the
accuracy requirements inspectable before optimization.

Evidence includes 306 generated hydrothermal dispatch models, matched kernel and Ruiz
controls, enumerated stress tests on three solvers, and independent batch scheduling.
A prospective 24-scenario study compares precision and complete-process costs
against unscaled and numerical controls. A larger fixed-policy confirmation
adds 3,840 assignments across 160 generated models and two solvers, with common
verified-candidate targets, whole-parent uncertainty and all adverse outcomes
retained. Together these studies locate the method's practical scope: explicit
precision specification and verifiable selection, with configuration-dependent
computational benefits and no general incremental time advantage over
conventional Euclidean scaling or certificate of mixed-integer optimality.

The combination of mathematical programming, computational evaluation and practical
guidance for model developers fits the journal's readership across operations research
methods and applications. Online Resource~1 supplies proofs and full protocols;
Online Resource~2 supplies code, data, saved outcomes and independent replay programs.
The historical public baseline is distinguished from the current local snapshot.

I am the sole author and declare no competing interests. No external financial
funding or research grant was received; ClusterUY computing resources
are acknowledged. The manuscript documents assistive AI use and my responsibility
for the scientific content. This work has not been published as a journal article
and is not under consideration elsewhere.

\bigskip
\noindent Sincerely,\\[6pt]
Mathias Rodr\'iguez Castro
\end{document}
